\paragraph{Coefficient-surface results.}
\begin{itemize}
\item \emph{Serial dependence reverses its effect with clearing efficacy.} At
$e=1$, raising $\rho$ from 0 to .25, .5 and .75 lowers the original-workflow
limit from 75.13\% to 56.73\%, 53.65\% and 52.64\%. At $e=.9$, .8 and .6 the
same change raises it (42.60$\to$49.60\%, 27.99$\to$42.64\% and
14.93$\to$27.06\%). With $\Delta(e)=F^*(\rho{=}.5,e)-F^*(\rho{=}0,e)$, $\Delta(1)=\mn21.5$,
$\Delta(.9)=+4.9$, $\Delta(.8)=+9.0$ and $\Delta(.6)=+5.6$ percentage points, so the sign
changes between $e=.9$ and $e=1$. The isolated two-opportunity bridge changes sign at
$e^*=.92353$ (Proposition~\ref{prop:interaction}). Because $\Delta(e)$ is
strongly nonlinear in $e$, a two-point interpolation between $e=.9$ and 1
understates the crossing; the finer original-workflow extension locates it
between .93 and .94 (.934; Section~\ref{sec:s-coef-extension}). The historical
prediction scores below, which used the protocol's two-point statistic, are
retained unchanged.
\item \emph{Deviations from the two-opportunity comparator.}
For independent misses with perfect clearing about 99\% of the baseline
coefficient comes from bridges with exactly two opportunities
(Table~\ref{tab:kmix}), so $F^*=75.13\%$ is essentially the analytic
$1-(.05/.10)^2=75.0\%$ plus a small correction from longer bridges. Structure
matters elsewhere: at $\rho=.5$, $e=1$, 26.2\% of the baseline coefficient comes
from bridges with four or more opportunities and the limit is 53.65\% against
52.27\% for the isolated bridge; at $\rho=.75$, $e=.6$ the limit is 27.06\%
against 16.18\%; for imperfect clearing it is 27.99\% against 26.53\%. The
opportunity-count mix is recorded only for the three primary profiles; for
other rows only the comparisons with $q_2$ and with the renewal approximation
below are available.
\item \emph{Renewal approximation (post hoc, round 1).} Suppose a contaminated
hand meets a susceptible patient after every second opportunity and the
contamination episode is unbounded and stationary. The expected number of
susceptible contacts per episode is then $S(h)=\sum_{j\ge1}q_{2j}(h)=\pi\Gamma
P\Gamma\{I-(P\Gamma)^2\}^{-1}\boldsymbol 1$, and the response is
$F_{\mathrm{geo}}=1-S(h_1)/S(h_0)$. It reproduces all 22 rows of the surface to
within 2.8 percentage points (13 rows within 0.5 and 18 within 1), against up to
10.9 points for the isolated $q_2$ comparator, and it exceeds the original-workflow
limit in every row, most at $\rho=.75$, $e=.6$ (29.84\% against 27.06\%;
Table~\ref{tab:surface-diff}). Its crossing at $\rho=.5$ is $e=.936$. On the
prespecified 132-row grid of Section~\ref{sec:s-coef-extension} it fails the
prespecified 3.5-point bound (4.32 points at $\rho=.95$, $e=.5$); post hoc, it
again exceeds $F^*$ in every row, by at most 2.48 points for $\rho\le.7$, a cut
chosen after the failure was seen. $S(h)$ is the
discrete-opportunity analogue of the worker contamination period in
Ross--Macdonald-type ward models \citep{austin1999}: under memoryless clearing it
depends only on $q_1=1-he$, whereas under a Markov law it is not determined by
$(h,q_1)$. The fit suggests that the ICU structure acts mainly through the
spacing of susceptible contacts (two opportunities) and the truncation of
episodes by resets; it is an approximation, not a theorem, and may partly reflect
the near-uniform raw spectrum over small even counts (Table~\ref{tab:raw-bridges}).
\item \emph{Completion range.} The ordering independent $>$ clustered $>$
imperfect holds at the three completion ranges computed in both workflows, with
changing gaps (Table~\ref{tab:compliance}), and, in the original-workflow limit,
at all five ranges of the extension, including 60\% to 70\% (47.45, 36.80 and
33.85\%) and 70\% to 80\% (57.65, 43.16 and 37.14\%;
Table~\ref{tab:extension-ranges}; post hoc, descriptive). On the isolated
two-opportunity bridge it fails at 40\% to 50\%, where imperfect clearing
narrowly exceeds clustering (Table~\ref{tab:estar}); no ICU limit was computed
below 60\%. With imperfect clearing,
raising completion from 95\% to 99\% has a limit of only 25.9\%. The
original-workflow extension tests the combined profile and lower ranges in
Section~\ref{sec:s-coef-extension}; it includes a failed directional prediction.
\end{itemize}

